% Tipo di documento. L'uso di twoside implica che i capitoli inizino sempre con la prima pagina a sinistra, eventualmente lasciando una pagina vuota nel capitolo precedente. Se questa cosa è fastidiosa, è possibile rimuoverlo. 
\documentclass[a4paper, twoside,openright]{report}

\usepackage[a4paper,top=3cm,bottom=3cm,left=3cm,right=3cm]{geometry} 
\usepackage[fontsize=11pt]{scrextend}
\usepackage[english]{babel}
\usepackage[fixlanguage]{babelbib}
\usepackage[utf8]{inputenc} 
\usepackage[T1]{fontenc}
\usepackage{fancyhdr}               
\usepackage{amssymb}
\usepackage{amsmath}
\usepackage{amsthm}         
\usepackage{graphicx}
\usepackage[dvipsnames]{xcolor}         
\usepackage{listings}          
\usepackage{hyperref}     
\usepackage[normalem]{ulem}

\pagestyle{fancy}
\fancyhf{}
\lhead{\rightmarki}
\rhead{\textbf{\thepage}}
\fancyfoot{}
\setlength{\headheight}{12.5pt}

% Rimuove il numero di pagina all'inizio dei capitoli
\fancypagestyle{plain}{
  \fancyfoot{}
  \fancyhead{}
  \renewcommand{\headrulewidth}{0pt}
}

% Stile del codice
\lstdefinestyle{codeStyle}{
    % Colore dei commenti
    commentstyle=\color{teal},
    % Colore delle keyword
    keywordstyle=\color{Magenta},
    % Stile dei numeri di riga
    numberstyle=\tiny\color{gray},
    % Colore delle stringhe
    stringstyle=\color{violet},
    % Dimensione e stile del testo
    basicstyle=\ttfamily\footnotesize,
    % newline solo ai whitespaces
    breakatwhitespace=false,     
    % newline si/no
    breaklines=true,                 
    % Posizione della caption, top/bottom 
    captionpos=b,                    
    % Mantiene gli spazi nel codice, utile per l'indentazione
    keepspaces=true,                 
    % Dove visualizzare i numeri di linea
    numbers=left,                    
    % Distanza tra i numeri di linea
    numbersep=5pt,                  
    % Mostra gli spazi bianchi o meno
    showspaces=false,                
    % Mostra gli spazi bianchi nelle stringhe
    showstringspaces=false,
    % Mostra i tab
    showtabs=false,
    % Dimensione dei tab
    tabsize=2
} \lstset{style=codeStyle}

\lstdefinestyle{longBlock}{
    commentstyle=\color{teal},
    keywordstyle=\color{Magenta},
    numberstyle=\tiny\color{gray},
    stringstyle=\color{violet},
    basicstyle=\ttfamily\scriptsize,
    breakatwhitespace=false,         
    breaklines=true,                 
    captionpos=b,                    
    keepspaces=true,                 
    numbers=left,                    
    numbersep=5pt,                  
    showspaces=false,                
    showstringspaces=false,
    showtabs=false,                  
    tabsize=2
} \lstset{style=codeStyle}

% Add padding between code section
\lstset{aboveskip=20pt,belowskip=20pt}

\hypersetup{
    colorlinks,
    linkcolor=CornflowerBlue,
    citecolor=CornflowerBlue
}

% Aggiunti definizioni, teoremi, linea e listing
\newtheorem{definition}{Definition}[section]
\newtheorem{theorem}{Theorem}[section]
\providecommand*\definitionautorefname{Definition}
\providecommand*\theoremautorefname{Theorem}
\providecommand*{\listingautorefname}{Listing}
\providecommand*\lstnumberautorefname{Line}

\raggedbottom

\begin{document}

\begin{titlepage}
\begin{figure}[!htb]
    \centering
    \includegraphics[keepaspectratio=true,scale=0.8]{format/img/LOGO_UNISI_VERTICALE_NERO_medio.jpg}
\end{figure}

\begin{center}
    \large{\bf DIPARTIMENTO DI SCIENZE MEDICHE CHIRURGICHE E NEUROSCIENZE
    DIPARTIMENTO DI MEDICINA MOLECOLARE E DELLO SVILUPPO\\
    DIPARTIMENTO DI BIOTECNOLOGIE MEDICHE}
    \vspace{0.5mm}
    \\ \large{CORSO DI LAUREA IN INFERMIERISTICA}
\end{center}

\vspace{15mm}
\begin{center}
    {\LARGE{\bf L'infermiere e il linfoma non-Hodgkin:\\ \vspace{5mm} dalle terapie tradizionali alle CAR-T }}
    
    % Se il titolo è abbastanza corto da stare su una riga, si può usare
    
    % {\LARGE{\bf Un fantastico titolo per la mia tesi!}}
\end{center}
\vspace{30mm}

\begin{minipage}[t]{0.47\textwidth}
	{\large{Relatore:}{\normalsize\vspace{3mm}
    \bf\\ \large{Prof: Linda Badii} \normalsize\vspace{3mm}\bf}}
\end{minipage}
\hfill
\begin{minipage}[t]{0.47\textwidth}\raggedleft
	{\large{Candidato:}{\normalsize\vspace{3mm} \bf\\ \large{Sara Durante}}}
\end{minipage}

\vspace{10mm}

\begin{center}
    {\large{Anno Accademico 2021-2022}}
\end{center}

\vspace{30mm}

\end{titlepage}


\tableofcontents

% Rimuovere se non si vuole la tabella delle figure
\listoffigures

\appendix

\bibliographystyle{plain}
\bibliography{bibliography.bib}

\end{document}